\begin{table}[htbp]\centering\small
\caption{Tres modelos de cielo claro sobre la misma geometría y el mismo estado atmosférico. La dispersión crece de un 26\,\% a masa de aire 4 hasta un factor 3,0 a masa de aire 11, y esa dispersión es la incertidumbre real de la cadena radiométrica, no el error formal de ninguno de ellos. El paper adopta SPECTRL2, que es el más alto de los tres, porque es la elección conservadora para el riesgo ocular.}
\label{tab:ensemble}
\begin{tabular}{lS[table-format=2.2]S[table-format=2.2]S[table-format=3.1]S[table-format=3.1]S[table-format=3.1]S[table-format=1.2]}
\toprule
Instante & {Altura (\textdegree)} & {Masa aire} & {SPECTRL2} & {Bird--Hulstrom} & {Ineichen--Perez} & {Razón máx/mín} \\
 & & & \multicolumn{3}{c}{DNI sin eclipse (\si{\watt\per\square\meter})} & \\
\midrule
C1 & 14.70 & 3.9 & 489.1 & 467.1 & 389.7 & 1.26 \\
30 min antes del máximo & 10.17 & 5.5 & 374.4 & 348.6 & 252.1 & 1.49 \\
instante de la totalidad & 4.75 & 10.7 & 185.5 & 154.0 & 61.1 & 3.04 \\
\bottomrule
\end{tabular}
\end{table}
